\documentclass[11pt,letterpaper]{article}

\usepackage[margin=1.4in]{geometry}
\usepackage{hyperref}
\usepackage{microtype}
\usepackage{parskip}

\hypersetup{
  colorlinks=true,
  linkcolor=black,
  urlcolor=black,
  pdftitle={Secret Sats},
  pdfauthor={Ross}
}

\title{Secret Sats}
\author{Ross \\ \small\href{mailto:ross@tacit.finance}{\texttt{ross@tacit.finance}}}
\date{September 2026}

\begin{document}
\maketitle
\thispagestyle{empty}

\begin{abstract}
\noindent
Secret Sats lets you pay in cBTC, a bitcoin-backed Tacit asset, on Bitcoin without revealing how much,
and inside a shielded pool without revealing who paid whom.
A self-custody lock on Bitcoin, backed by an Ethereum escrow, mints cBTC in Tacit's confidential pool,
and an SP1 light-client proof of Ethereum, not a bridge multisig, carries it back to Bitcoin as a native
Tacit asset.
There, every transfer hides its amount, and a transfer to a stealth address pays a one-time key that only
the recipient can find.
Shielding it into Tacit's Bitcoin pool also hides where a payment came from, and relayers can post pool
spends so the sender needs no Bitcoin wallet.
Sats are meant to enter and leave by atomic swap, and by covenant once Bitcoin allows one; the cBTC pilot
is live on mainnet, and the pool's shield, private payment and exit have run end to end on signet.
\end{abstract}

\vfill
\begin{center}
\small\url{https://github.com/z0r0z/tacit}
\end{center}

\end{document}
